%***************************************************************************************************
% Podstawowe ustawienia języka, według którego formatowany będzie dokument
\usepackage[polish]{babel}
\usepackage{diagbox}

% Pakiet babel dla polskiego języka powoduje konflikt z pakietem amssymb.
\let\lll\undefined

% W przypadku wielojęzykowości ustawia główny język dokumentu
\selectlanguage{polish}

% Kodowanie dokumentu
\usepackage[utf8]{inputenc} 
\usepackage[T1]{fontenc} 

% Dowolny rozmiar czcionek, kodowanie znaków
\usepackage{lmodern}

% Polskie wcięcia akapitów
\usepackage{indentfirst}

% Polskie formatowanie typograficzne
\frenchspacing

% Zapewnia liczne usprawnienia wyświetlania i organizacji matematycznych formuł. 
\usepackage{amsmath}
\usepackage{amssymb}
\usepackage{mathrsfs}
\usepackage{mathtools}

% ***************************************************************************************************
% Kolory  
% ***************************************************************************************************
% Umożliwia kolorowanie poszczególnych komórek tabeli i rozszerzoną kontrolę
\usepackage[table, svgnames]{xcolor}

% Umożliwia łatwą zmianę koloru linii w tabeli
\usepackage{tabu}

% Definicje kolorów
\definecolor{lgray}{HTML}{9F9F9F}
\definecolor{dgray}{HTML}{5F5F5F}

% ***************************************************************************************************
% Algorytmy i Listingi
% ***************************************************************************************************
\usepackage[ruled,vlined,linesnumbered,longend,algochapter]{algorithm2e}

% Zamiana nazwy środowiska z domyślnej "Algorithm X" na "Pseudokod X"
\newenvironment{algorithm-custom}[1][htb]{
	\renewcommand{\algorithmcfname}{Pseudokod}
	\begin{algorithm}[#1]%
	}{
\end{algorithm}
}

\newcommand\algcomment[1]{\footnotesize{#1}}
\SetCommentSty{algcomment}

% Wyśrodkowana tylda
\usepackage{textcomp}
\newcommand{\textapprox}{\raisebox{0.5ex}{\texttildelow}}

% Listowanie algorytmów
\renewcommand\algorithmcfname{Algorytm}
\renewcommand\listalgorithmcfname{Spis algorytmów}

% Pakiet listings i pełna polska konfiguracja (kodowanie UTF-8 oraz polskie nazwy spisów)
\usepackage{listings}
\renewcommand{\lstlistingname}{Listing}
\renewcommand{\lstlistlistingname}{Spis listingów}

\addto\captionspolish{%
	\renewcommand{\lstlistingname}{Listing}%
	\renewcommand{\lstlistlistingname}{Spis listingów}%
}

\lstset{
	inputencoding=utf8,
	extendedchars=true,
	literate=
		{ą}{{\k{a}}}1 {ć}{{\'c}}1 {ę}{{\k{e}}}1 {ł}{{\l{}}}1 {ń}{{\'n}}1
		{ó}{{\'o}}1 {ś}{{\'s}}1 {ź}{{\'z}}1 {ż}{{\.z}}1
		{Ą}{{\k{A}}}1 {Ć}{{\'C}}1 {Ę}{{\k{E}}}1 {Ł}{{\L{}}}1 {Ń}{{\'N}}1
		{Ó}{{\'O}}1 {Ś}{{\'S}}1 {Ź}{{\'Z}}1 {Ż}{{\.Z}}1
}

% ***************************************************************************************************
% Marginesy 
% ***************************************************************************************************
\usepackage[headheight=18pt, top=25mm, bottom=25mm, left=25mm, right=25mm]{geometry}

\makeatletter
\setlength\@fptop{0\p@}	
\makeatother

% ***************************************************************************************************
% Styl 
% ***************************************************************************************************
\usepackage{titling}
\usepackage[subfigure]{tocloft}	
\tocloftpagestyle{tableOfContentStyle}

\usepackage{fancyhdr}

% Domyślny czysty styl (używany w frontmatter - bez przedmiotu, bez numeru)
\fancypagestyle{plain}{
	\fancyhf{}
	\renewcommand{\headrulewidth}{0pt}
	\renewcommand{\footrulewidth}{0pt}
}

% Domyślny styl dla stron właściwych (środek rozdziału)
\fancypagestyle{custom}{
	\fancyhf{}
	\fancyhead[RO]{
		\hrulefill \hspace{16pt} \large Rozdział \thechapter
		\put(-471.5,5.5){\makebox(0,0)[l]{\small Politechnika Wrocławska}}
	}
	\fancyhead[LE]{
		\large Rozdział \thechapter \hspace{16pt} \hrulefill 
		\put(-188,5.5){\makebox(0,0)[l]{\small Wydział Informatyki i Telekomunikacji}}
	}
	\fancyfoot[LE,RO]{\thepage}
	\fancyfoot[LO,RE]{\small\Przedmiot}
	\renewcommand{\headrulewidth}{0pt}
	\renewcommand{\footrulewidth}{0.2pt}
}

% Styl dla List of Figures (BEZ PRZEDMIOTU)
\fancypagestyle{ListofFiguresStyle}{
	\fancyhf{}
	\fancyhead[RO]{
		\hrulefill \hspace{16pt} \large Spis rysunków
		\put(-471.5,5.5){\makebox(0,0)[l]{\small Politechnika Wrocławska}}
	}
	\fancyhead[LE]{
		\large Spis rysunków \hspace{16pt} \hrulefill 
		\put(-188,5.5){\makebox(0,0)[l]{\small Wydział Informatyki i Telekomunikacji}}
	}
	\fancyfoot[LE,RO]{\thepage}
	\renewcommand{\headrulewidth}{0pt}
	\renewcommand{\footrulewidth}{0.2pt}
}

% Styl dla List of Tables (BEZ PRZEDMIOTU)
\fancypagestyle{ListofTablesStyle}{
	\fancyhf{}
	\fancyhead[RO]{
		\hrulefill \hspace{16pt} \large Spis tabel
		\put(-471.5,5.5){\makebox(0,0)[l]{\small Politechnika Wrocławska}}
	}
	\fancyhead[LE]{
		\large Spis tabel \hspace{16pt} \hrulefill 
		\put(-188,5.5){\makebox(0,0)[l]{\small Wydział Informatyki i Telekomunikacji}}
	}
	\fancyfoot[LE,RO]{\thepage}
	\renewcommand{\headrulewidth}{0pt}
	\renewcommand{\footrulewidth}{0.2pt}
}

% Styl dla List of Listings
\fancypagestyle{ListofListingsStyle}{
	\fancyhf{}
	\fancyhead[RO]{
		\hrulefill \hspace{16pt} \large Spis listingów
		\put(-471.5,5.5){\makebox(0,0)[l]{\small Politechnika Wrocławska}}
	}
	\fancyhead[LE]{
		\large Spis listingów \hspace{16pt} \hrulefill 
		\put(-188,5.5){\makebox(0,0)[l]{\small Wydział Informatyki i Telekomunikacji}}
	}
	\fancyfoot[LE,RO]{\thepage}
	\renewcommand{\headrulewidth}{0pt}
	\renewcommand{\footrulewidth}{0.2pt}
}

% Styl dla List of Algorithms
\fancypagestyle{ListofAlgorithmsStyle}{
	\fancyhf{}
	\fancyhead[RO]{
		\hrulefill \hspace{16pt} \large Spis algorytmów
		\put(-471.5,5.5){\makebox(0,0)[l]{\small Politechnika Wrocławska}}
	}
	\fancyhead[LE]{
		\large Spis algorytmów \hspace{16pt} \hrulefill 
		\put(-188,5.5){\makebox(0,0)[l]{\small Wydział Informatyki i Telekomunikacji}}
	}
	\fancyfoot[LE,RO]{\thepage}
	\renewcommand{\headrulewidth}{0pt}
	\renewcommand{\footrulewidth}{0.2pt}
}

% Styl dla bibliografii (MA PRZEDMIOT)
\fancypagestyle{bibliographyStyle}{
	\fancyhf{}
	\fancyhead[RO]{
		\hrulefill \hspace{16pt} \large Bibliografia
		\put(-471.5,5.5){\makebox(0,0)[l]{\small Politechnika Wrocławska}}
	}
	\fancyhead[LE]{
		\large Bibliografia \hspace{16pt} \hrulefill 
		\put(-188,5.5){\makebox(0,0)[l]{\small Wydział Informatyki i Telekomunikacji}}
	}
	\fancyfoot[LE,RO]{\thepage}
	\fancyfoot[LO,RE]{\small\Przedmiot}
	\renewcommand{\headrulewidth}{0pt}
	\renewcommand{\footrulewidth}{0.2pt}
}

% Styl dla dodatków (MA PRZEDMIOT)
\fancypagestyle{appendixStyle}{
	\fancyhf{}
	\fancyhead[RO]{
		\hrulefill \hspace{16pt} \large Dodatki \thechapter
		\put(-472.1, 12.1){\makebox(0,0)[l]{\includegraphics[width=0.05\textwidth]{resources/pwr-logo}}}
		\put(-443,5.5){\makebox(0,0)[l]{\small Politechnika Wrocławska}}
	}
	\fancyhead[LE]{
		\large Dodatek \thechapter \hspace{16pt} \hrulefill 
		\put(-22, 12.1){\makebox(0,0)[l]{\includegraphics[width=0.05\textwidth]{wiz-logo}}}
		\put(-141,5.5){\makebox(0,0)[l]{\small Wydział Informatyki i Telekomunikacji}}
	}
	\fancyfoot[LE,RO]{\thepage}
	\fancyfoot[LO,RE]{\small\Przedmiot}
	\renewcommand{\headrulewidth}{0pt}
	\renewcommand{\footrulewidth}{0.2pt}
}

% Styl dla List of Abbreviation
\fancypagestyle{ListofAbbreviationStyle}{
	\fancyhf{}
	\fancyhead[RO]{
		\hrulefill \hspace{16pt} \large Spis skrótów
		\put(-471.5,5.5){\makebox(0,0)[l]{\small Politechnika Wrocławska}}
	}
	\fancyhead[LE]{
		\large Spis skrótów \hspace{16pt} \hrulefill 
		\put(-188,5.5){\makebox(0,0)[l]{\small Wydział Informatyki i Telekomunikacji}}
	}
	\fancyfoot[LE,RO]{\thepage}
	\renewcommand{\headrulewidth}{0pt}
	\renewcommand{\footrulewidth}{0.2pt}
}

% Styl dla pozostałych stron spisu treści (czysty)
\fancypagestyle{tableOfContentStyle}{
	\fancyhf{}
	\renewcommand{\headrulewidth}{0pt}
	\renewcommand{\footrulewidth}{0.2pt}
}

% Formatowanie tytułów rozdziałów i sekcji (wraz z subsubsection)
\usepackage{titlesec}
\setcounter{secnumdepth}{3}
\setcounter{tocdepth}{2}

\titleformat{\chapter}[hang]
{\vspace{-10ex}\normalfont\Huge\bfseries}{\thechapter.}{1ex}{} 

\titleformat{\section}[hang]
{\normalfont\Large\bfseries}{\thesection.}{1ex}{} 

\titleformat{\subsection}[hang]
{\normalfont\large\bfseries}{\thesubsection.}{1ex}{} 

\titleformat{\subsubsection}[hang]
{\normalfont\normalsize\bfseries}{\thesubsubsection.}{1ex}{}

% ***************************************************************************************************
% Linki
% ***************************************************************************************************
\usepackage{hyperref}
\hypersetup{
	colorlinks	=	true,
	linkcolor	=	blue,
	citecolor	=	blue,
	urlcolor	=	blue
}
\urlstyle{same}

% ***************************************************************************************************
% Listy
% ***************************************************************************************************
\usepackage{enumitem}
\newlist{myitemize}{itemize}{3}
\setlist[myitemize,1]{label=\textbullet, leftmargin=4mm}
\setlist[myitemize,2]{label=$\diamond$, leftmargin=8mm}
\setlist[myitemize,3]{label=$\diamond$, leftmargin=12mm}

% ***************************************************************************************************
% Inne pakiety
% ***************************************************************************************************
\usepackage{graphicx}
\usepackage{float}
\usepackage{caption}
\usepackage{footnote}
\usepackage{dirtree}
\usepackage{multirow}
\usepackage{calc}

% ***************************************************************************************************
% Matematyczne skróty
% ***************************************************************************************************
\newcommand{\RR}{\mathbb{R}}
\newcommand{\NN}{\mathbb{N}}
\newcommand{\QQ}{\mathbb{Q}}
\newcommand{\ZZ}{\mathbb{Z}}
\newcommand{\IMP}{\rightarrow}
\newcommand{\IFF}{\leftrightarrow}

% ***************************************************************************************************
% Środowiska
% ***************************************************************************************************
\newtheorem{theorem}{Twierdzenie}[chapter]
\newtheorem{lemma}{Lemat}[chapter]
\newtheorem{example}{Przykład}[chapter]
\newtheorem{corollary}{Wniosek}[chapter]
\newtheorem{definition}{Definicja}[chapter]

\newenvironment{proof}{\par\noindent \textbf{Dowód.}}{\begin{flushright}\vspace*{-6mm}\mbox{$\blacklozenge$}\end{flushright}}
\newenvironment{remark}{\bigskip \par\noindent \small \textbf{Uwaga.}}{\begin{small}\vspace*{4mm}\end{small}}

% Dodatkowe pakiety pomocnicze
\usepackage{psfrag}
\usepackage{amsfonts}
\usepackage{supertabular}
\usepackage{array}
\usepackage{tabularx}
\usepackage{hhline}
\usepackage{microtype}
\usepackage{booktabs}
\usepackage{makecell}
\usepackage{rotating}
\usepackage{multicol}
\usepackage{cuted}
\usepackage{colortbl}
\usepackage{adjustbox}
\usepackage{color,soul}
\usepackage{subfigure}
\usepackage{pdfpages}

% Wyłączenie generowania pustych stron (zastąpienie cleardoublepage zwykłym clearpage)
\renewcommand{\cleardoublepage}{\clearpage}

\usepackage{pifont}
\newcommand{\cmark}{\ding{51}}
\newcommand{\xmark}{\ding{55}}
\newcommand{\bftab}{\fontseries{b}\selectfont}

\newcolumntype{P}[1]{>{\raggedright\arraybackslash\noindent}p{#1}}
\newcolumntype{R}[1]{>{\raggedleft\arraybackslash}p{#1}}
\newcolumntype{L}[1]{>{\raggedright\arraybackslash}p{#1}}
\newcolumntype{C}[1]{>{\centering\arraybackslash}m{#1}}

\usepackage[nocompress]{cite}
\usepackage{svg}
\usepackage{tabto}
\usepackage{wrapfig}